%-----------------------------------------------------------------------------------------------------------------------------------------------%
%	The MIT License (MIT)
%
%	Copyright (c) 2021 Jitin Nair
%
%	Permission is hereby granted, free of charge, to any person obtaining a copy
%	of this software and associated documentation files (the "Software"), to deal
%	in the Software without restriction, including without limitation the rights
%	to use, copy, modify, merge, publish, distribute, sublicense, and/or sell
%	copies of the Software, and to permit persons to whom the Software is
%	furnished to do so, subject to the following conditions:
%	
%	THE SOFTWARE IS PROVIDED "AS IS", WITHOUT WARRANTY OF ANY KIND, EXPRESS OR
%	IMPLIED, INCLUDING BUT NOT LIMITED TO THE WARRANTIES OF MERCHANTABILITY,
%	FITNESS FOR A PARTICULAR PURPOSE AND NONINFRINGEMENT. IN NO EVENT SHALL THE
%	AUTHORS OR COPYRIGHT HOLDERS BE LIABLE FOR ANY CLAIM, DAMAGES OR OTHER
%	LIABILITY, WHETHER IN AN ACTION OF CONTRACT, TORT OR OTHERWISE, ARISING FROM,
%	OUT OF OR IN CONNECTION WITH THE SOFTWARE OR THE USE OR OTHER DEALINGS IN
%	THE SOFTWARE.
%	
%
%-----------------------------------------------------------------------------------------------------------------------------------------------%

%----------------------------------------------------------------------------------------
%	DOCUMENT DEFINITION
%----------------------------------------------------------------------------------------

% article class because we want to fully customize the page and not use a cv template
\documentclass[letterpaper,10pt]{article}

% Adding packages for skills
% -------------------------
\usepackage{geometry}
\geometry{margin=0.75in}

\usepackage{array}
\usepackage[svgnames]{xcolor}
\usepackage{tikz}
\usetikzlibrary{math}
\usepackage{progressbar}
%\usepackage[maxbibnames=99]{biblatex}
\usepackage{url}
%\usepackage[backend=biber,style=numeric,sorting=ydnt,defernumbers=true]{biblatex}
% Bibliography is typeset manually in the Publications / Talks sections below
% (rendered from publications.md); biblatex/biber is intentionally not used.
% \usepackage[style=numeric, sorting=none, backend=biber, maxbibnames=99]{biblatex}
% \addbibresource{citations.bib}


%----------------------------------------------------------------------------------------
%	FONT
%----------------------------------------------------------------------------------------

% Same font stack as the research and teaching statements:
% lmodern first gives T1 a scalable fallback so nothing lands on a bitmap
% Computer Modern face; 'default' is what actually makes Lato the document
% font (without it the package loads but is never used).
\usepackage{lmodern}
\usepackage[T1]{fontenc}
\usepackage[scaled=0.98,default]{lato}


% % fontspec allows you to use TTF/OTF fonts directly
% \usepackage{fontspec}
% \defaultfontfeatures{Ligatures=TeX}

% % modified for ShareLaTeX use
% \setmainfont[
% SmallCapsFont = Fontin-SmallCaps.otf,
% BoldFont = Fontin-Bold.otf,
% ItalicFont = Fontin-Italic.otf
% ]
% {Fontin.otf}

%----------------------------------------------------------------------------------------
%	PACKAGES
%----------------------------------------------------------------------------------------
\usepackage{url}
\usepackage{parskip}
\setlength{\parskip}{3.5pt plus 1pt}   % tighter inter-paragraph spacing than parskip default

%other packages for formatting
\RequirePackage{color}
\RequirePackage{graphicx}
\usepackage[usenames,dvipsnames]{xcolor}
%tabularx environment
%for lists within experience section
\usepackage{enumitem}

% centered version of 'X' col. type
\newcolumntype{C}{>{\centering\arraybackslash}X} 

%to prevent spillover of tabular into next pages
\usepackage{supertabular}
\usepackage{tabularx}
\usepackage{needspace}   % reserve vertical space so a block isn't split across pages
\newlength{\fullcollw}
\setlength{\fullcollw}{0.47\textwidth}
\newlength{\awardcolw}   % description column of the page-breakable awards table

%custom \section
\usepackage{titlesec}				
\usepackage{multicol}
\usepackage{multirow}

%CV Sections inspired by: 
%http://stefano.italians.nl/archives/26
\titleformat{\section}{\normalsize\scshape\raggedright}{}{0em}{}[\titlerule]
%\titleformat{\section}{\textbf{}\normalsize\scshape\raggedright}{}{0em}{}[\titlerule]
\titlespacing{\section}{0pt}{7pt}{5pt}

%for publications

%Setup hyperref package, and colours for links
\definecolor{customcolor}{HTML}{6f784d}
\colorlet{color1}{customcolor}
\usepackage[unicode, draft=false]{hyperref}
\definecolor{linkcolour}{rgb}{0,0.2,0.6}
\definecolor{darknavy}{HTML}{1F3864}   % dark navy for hyperlinks (prints legibly)
\hypersetup{colorlinks,breaklinks,urlcolor=darknavy,linkcolor=darknavy,citecolor=darknavy}
\usepackage{xurl}   % allow long DOIs/URLs to wrap across lines
\urlstyle{same}     % use the body font for wrapped identifiers, not typewriter
\hypersetup{pdftitle={Jenna M. Kline --- Curriculum Vitae}, pdfauthor={Jenna M. Kline}}

%--- Publication-list formatting (manual, rendered from publications.md) ---
\providecommand{\emailsymbol}{\faEnvelope~}
\providecommand{\homepagesymbol}{\faGlobe~}
% Compact small-caps subsection headers for publication categories
\titleformat{\subsection}{\normalsize\bfseries\raggedright}{}{0em}{}
\titlespacing{\subsection}{0pt}{5pt}{2pt}
% Hanging-indent list for references. Numbering restarts in every category and
% carries a category prefix (F/J/C/P/T/A) so cross-references stay unambiguous.
\newlist{pubs}{enumerate}{1}
\setlist[pubs]{leftmargin=2.5em, labelindent=0pt, labelwidth=2.1em, labelsep=0.4em, align=left, label=\arabic*., ref=\arabic*, itemsep=0.45em, topsep=0.12em, parsep=0pt}
% Year sub-heading within a category
\newcommand{\pubyear}[1]{\smallskip\noindent\textbf{#1}\par\nobreak\vspace{0.15em}}
% Clickable identifiers
\newcommand{\doilink}[1]{\href{https://doi.org/#1}{doi:\nolinkurl{#1}}}
% Artifact host link (GitHub, Hugging Face, ...). \textnormal keeps it upright
% when it sits inside the italic venue field of \pvenue.
\newcommand{\phost}[2]{\textnormal{\href{#1}{#2}}}

% --- Three-line publication entry: title / authors / venue + year -----------
% Title carries the emphasis; venue and year sit on their own line so a reader
% can scan where and when the work appeared without parsing the author list.
\newcommand{\ptitle}[1]{\textbf{#1}\par\nobreak\vspace{0.4pt}}
\newcommand{\pauth}[1]{#1\par\nobreak\vspace{0.4pt}}
% #1 = venue, #2 = year, #3 = identifiers, awards, and status notes (may be empty)
\newcommand{\pvenue}[3]{\textit{#1}, \textbf{#2}. #3\par}
% Artifact detail line: everything but the title and the download count runs on
% as one paragraph. #1 = kind, hosts, and role (upright); #2 = contents and the
% publications it supports (italic, like the venue field of \pvenue).
\newcommand{\pdetail}[2]{#1 \textit{#2}\par}
% Status flag for work that is not yet published
\newcommand{\status}[1]{\textbf{[#1]}}
% Category sub-heading inside a publication list block
\newcommand{\talkgroup}[1]{\smallskip\noindent\textit{\textbf{#1}}\par\nobreak\vspace{0.2em}}
% Impact metric on its own line, below the venue line of an artifact entry
\newcommand{\pmetric}[1]{\textbf{#1}\par}
% Qualifier trailing a subsection heading on the same line. \normalfont undoes
% the heading's bold; \footnotesize\itshape keeps it subordinate to the title.
\newcommand{\subnote}[1]{{\normalfont\footnotesize\itshape #1}}
% Field-experience entry: location (left) + dates (right) on line 1, description on line 2
% #1 = location, #2 = dates, #3 = description
\newcommand{\fieldentry}[3]{%
  \noindent\begin{tabularx}{\linewidth}{@{}X r@{}}\textbf{#1} & #2\end{tabularx}\par\nobreak\vspace{1pt}%
  {\noindent #3\par}\vspace{2pt}}

%for social icons
\usepackage{fontawesome5}
\usepackage{fontawesome5}
\usepackage{graphicx}  

% for bib
%\usepackage{cite}



% %for header
% \usepackage{fancyhdr}
% \pagestyle{fancy}
% \fancyhf{}
% \fancyfoot[R]{Page \thepage \hspace{1pt} of 4}
% \fancyfoot[L]{Jenna M. Kline}
% \renewcommand{\headrulewidth}{0pt}
% \renewcommand{\footrulewidth}{0pt}

%----------------------------------------------------------------------------------------
%	HEADER AND FOOTER DEFINITIONS
%----------------------------------------------------------------------------------------
\newcommand{\yourname}{Jenna M. Kline}
% \newcommand{\youremail}{jennamkline@gmail.com}
% \newcommand{\yourweb}{jennamk14.github.io}

\newcommand{\youremail}{\href{mailto:jennamk9@mit.edu}{jennamk9@mit.edu}}
\newcommand{\yourweb}{\href{https://jennamk14.github.io/}{jennamk14.github.io}}

%for header
\usepackage{fancyhdr}
\pagestyle{fancy}
\fancyhf{}
% \fancyhead[L]{\yourname\quad}
% \fancyhead[R]{Curriculum Vitae}
%\fancyfoot[R]{Page \thepage \hspace{1pt} of \pageref{LastPage}}
\fancyfoot[R]{Page \thepage \hspace{1pt}  of \pageref{LastPage}}
% \fancyfoot[C]{
%     \footnotesize
%     \textcolor{customcolor}{\itshape\yourweb}\quad
%     \textcolor{customcolor}{\youremail}}
\fancyfoot[C]{%
  \footnotesize
  \emailsymbol\youremail \quad $\bullet$ \quad
 \itshape\homepagesymbol\yourweb
}
    
\renewcommand{\headrulewidth}{0pt}
\renewcommand{\footrulewidth}{0pt}

% First page carries the contact block in the title, so suppress the
% email/website footer there (keep only the page number).
\fancypagestyle{firstpage}{%
  \fancyhf{}%
  \fancyfoot[R]{Page \thepage \hspace{1pt}  of \pageref{LastPage}}%
  \renewcommand{\headrulewidth}{0pt}%
  \renewcommand{\footrulewidth}{0pt}%
}

\usepackage{lastpage}

%\lhead{This is my name}
%\rhead{this is page \thepage}
%\renewcommand{\headrulewidth}{0.4pt}
%\renewcommand{\footrulewidth}{0.4pt}

%debug page outer frames
%\usepackage{showframe}

%-----------------------
%  commands for circles 
%----------------------------------------------------------------------------------------
%	BEGIN DOCUMENT
%----------------------------------------------------------------------------------------
%\pagenumbering{arabic}

% ---------------------------------------------------------------------------
% Per-application tailoring. Wrapper files (e.g., Kline_CV_washu.tex) define
% these BEFORE %-----------------------------------------------------------------------------------------------------------------------------------------------%
%	The MIT License (MIT)
%
%	Copyright (c) 2021 Jitin Nair
%
%	Permission is hereby granted, free of charge, to any person obtaining a copy
%	of this software and associated documentation files (the "Software"), to deal
%	in the Software without restriction, including without limitation the rights
%	to use, copy, modify, merge, publish, distribute, sublicense, and/or sell
%	copies of the Software, and to permit persons to whom the Software is
%	furnished to do so, subject to the following conditions:
%	
%	THE SOFTWARE IS PROVIDED "AS IS", WITHOUT WARRANTY OF ANY KIND, EXPRESS OR
%	IMPLIED, INCLUDING BUT NOT LIMITED TO THE WARRANTIES OF MERCHANTABILITY,
%	FITNESS FOR A PARTICULAR PURPOSE AND NONINFRINGEMENT. IN NO EVENT SHALL THE
%	AUTHORS OR COPYRIGHT HOLDERS BE LIABLE FOR ANY CLAIM, DAMAGES OR OTHER
%	LIABILITY, WHETHER IN AN ACTION OF CONTRACT, TORT OR OTHERWISE, ARISING FROM,
%	OUT OF OR IN CONNECTION WITH THE SOFTWARE OR THE USE OR OTHER DEALINGS IN
%	THE SOFTWARE.
%	
%
%-----------------------------------------------------------------------------------------------------------------------------------------------%

%----------------------------------------------------------------------------------------
%	DOCUMENT DEFINITION
%----------------------------------------------------------------------------------------

% article class because we want to fully customize the page and not use a cv template
\documentclass[letterpaper,10pt]{article}

% Adding packages for skills
% -------------------------
\usepackage{geometry}
\geometry{margin=0.75in}

\usepackage{array}
\usepackage[svgnames]{xcolor}
\usepackage{tikz}
\usetikzlibrary{math}
\usepackage{progressbar}
%\usepackage[maxbibnames=99]{biblatex}
\usepackage{url}
%\usepackage[backend=biber,style=numeric,sorting=ydnt,defernumbers=true]{biblatex}
% Bibliography is typeset manually in the Publications / Talks sections below
% (rendered from publications.md); biblatex/biber is intentionally not used.
% \usepackage[style=numeric, sorting=none, backend=biber, maxbibnames=99]{biblatex}
% \addbibresource{citations.bib}


%----------------------------------------------------------------------------------------
%	FONT
%----------------------------------------------------------------------------------------

% Same font stack as the research and teaching statements:
% lmodern first gives T1 a scalable fallback so nothing lands on a bitmap
% Computer Modern face; 'default' is what actually makes Lato the document
% font (without it the package loads but is never used).
\usepackage{lmodern}
\usepackage[T1]{fontenc}
\usepackage[scaled=0.98,default]{lato}


% % fontspec allows you to use TTF/OTF fonts directly
% \usepackage{fontspec}
% \defaultfontfeatures{Ligatures=TeX}

% % modified for ShareLaTeX use
% \setmainfont[
% SmallCapsFont = Fontin-SmallCaps.otf,
% BoldFont = Fontin-Bold.otf,
% ItalicFont = Fontin-Italic.otf
% ]
% {Fontin.otf}

%----------------------------------------------------------------------------------------
%	PACKAGES
%----------------------------------------------------------------------------------------
\usepackage{url}
\usepackage{parskip}
\setlength{\parskip}{3.5pt plus 1pt}   % tighter inter-paragraph spacing than parskip default

%other packages for formatting
\RequirePackage{color}
\RequirePackage{graphicx}
\usepackage[usenames,dvipsnames]{xcolor}
%tabularx environment
%for lists within experience section
\usepackage{enumitem}

% centered version of 'X' col. type
\newcolumntype{C}{>{\centering\arraybackslash}X} 

%to prevent spillover of tabular into next pages
\usepackage{supertabular}
\usepackage{tabularx}
\usepackage{needspace}   % reserve vertical space so a block isn't split across pages
\newlength{\fullcollw}
\setlength{\fullcollw}{0.47\textwidth}
\newlength{\awardcolw}   % description column of the page-breakable awards table

%custom \section
\usepackage{titlesec}				
\usepackage{multicol}
\usepackage{multirow}

%CV Sections inspired by: 
%http://stefano.italians.nl/archives/26
\titleformat{\section}{\normalsize\scshape\raggedright}{}{0em}{}[\titlerule]
%\titleformat{\section}{\textbf{}\normalsize\scshape\raggedright}{}{0em}{}[\titlerule]
\titlespacing{\section}{0pt}{7pt}{5pt}

%for publications

%Setup hyperref package, and colours for links
\definecolor{customcolor}{HTML}{6f784d}
\colorlet{color1}{customcolor}
\usepackage[unicode, draft=false]{hyperref}
\definecolor{linkcolour}{rgb}{0,0.2,0.6}
\definecolor{darknavy}{HTML}{1F3864}   % dark navy for hyperlinks (prints legibly)
\hypersetup{colorlinks,breaklinks,urlcolor=darknavy,linkcolor=darknavy,citecolor=darknavy}
\usepackage{xurl}   % allow long DOIs/URLs to wrap across lines
\urlstyle{same}     % use the body font for wrapped identifiers, not typewriter
\hypersetup{pdftitle={Jenna M. Kline --- Curriculum Vitae}, pdfauthor={Jenna M. Kline}}

%--- Publication-list formatting (manual, rendered from publications.md) ---
\providecommand{\emailsymbol}{\faEnvelope~}
\providecommand{\homepagesymbol}{\faGlobe~}
% Compact small-caps subsection headers for publication categories
\titleformat{\subsection}{\normalsize\bfseries\raggedright}{}{0em}{}
\titlespacing{\subsection}{0pt}{5pt}{2pt}
% Hanging-indent list for references. Numbering restarts in every category and
% carries a category prefix (F/J/C/P/T/A) so cross-references stay unambiguous.
\newlist{pubs}{enumerate}{1}
\setlist[pubs]{leftmargin=2.5em, labelindent=0pt, labelwidth=2.1em, labelsep=0.4em, align=left, label=\arabic*., ref=\arabic*, itemsep=0.45em, topsep=0.12em, parsep=0pt}
% Year sub-heading within a category
\newcommand{\pubyear}[1]{\smallskip\noindent\textbf{#1}\par\nobreak\vspace{0.15em}}
% Clickable identifiers
\newcommand{\doilink}[1]{\href{https://doi.org/#1}{doi:\nolinkurl{#1}}}
% Artifact host link (GitHub, Hugging Face, ...). \textnormal keeps it upright
% when it sits inside the italic venue field of \pvenue.
\newcommand{\phost}[2]{\textnormal{\href{#1}{#2}}}

% --- Three-line publication entry: title / authors / venue + year -----------
% Title carries the emphasis; venue and year sit on their own line so a reader
% can scan where and when the work appeared without parsing the author list.
\newcommand{\ptitle}[1]{\textbf{#1}\par\nobreak\vspace{0.4pt}}
\newcommand{\pauth}[1]{#1\par\nobreak\vspace{0.4pt}}
% #1 = venue, #2 = year, #3 = identifiers, awards, and status notes (may be empty)
\newcommand{\pvenue}[3]{\textit{#1}, \textbf{#2}. #3\par}
% Artifact detail line: everything but the title and the download count runs on
% as one paragraph. #1 = kind, hosts, and role (upright); #2 = contents and the
% publications it supports (italic, like the venue field of \pvenue).
\newcommand{\pdetail}[2]{#1 \textit{#2}\par}
% Status flag for work that is not yet published
\newcommand{\status}[1]{\textbf{[#1]}}
% Category sub-heading inside a publication list block
\newcommand{\talkgroup}[1]{\smallskip\noindent\textit{\textbf{#1}}\par\nobreak\vspace{0.2em}}
% Impact metric on its own line, below the venue line of an artifact entry
\newcommand{\pmetric}[1]{\textbf{#1}\par}
% Qualifier trailing a subsection heading on the same line. \normalfont undoes
% the heading's bold; \footnotesize\itshape keeps it subordinate to the title.
\newcommand{\subnote}[1]{{\normalfont\footnotesize\itshape #1}}
% Field-experience entry: location (left) + dates (right) on line 1, description on line 2
% #1 = location, #2 = dates, #3 = description
\newcommand{\fieldentry}[3]{%
  \noindent\begin{tabularx}{\linewidth}{@{}X r@{}}\textbf{#1} & #2\end{tabularx}\par\nobreak\vspace{1pt}%
  {\noindent #3\par}\vspace{2pt}}

%for social icons
\usepackage{fontawesome5}
\usepackage{fontawesome5}
\usepackage{graphicx}  

% for bib
%\usepackage{cite}



% %for header
% \usepackage{fancyhdr}
% \pagestyle{fancy}
% \fancyhf{}
% \fancyfoot[R]{Page \thepage \hspace{1pt} of 4}
% \fancyfoot[L]{Jenna M. Kline}
% \renewcommand{\headrulewidth}{0pt}
% \renewcommand{\footrulewidth}{0pt}

%----------------------------------------------------------------------------------------
%	HEADER AND FOOTER DEFINITIONS
%----------------------------------------------------------------------------------------
\newcommand{\yourname}{Jenna M. Kline}
% \newcommand{\youremail}{jennamkline@gmail.com}
% \newcommand{\yourweb}{jennamk14.github.io}

\newcommand{\youremail}{\href{mailto:jennamk9@mit.edu}{jennamk9@mit.edu}}
\newcommand{\yourweb}{\href{https://jennamk14.github.io/}{jennamk14.github.io}}

%for header
\usepackage{fancyhdr}
\pagestyle{fancy}
\fancyhf{}
% \fancyhead[L]{\yourname\quad}
% \fancyhead[R]{Curriculum Vitae}
%\fancyfoot[R]{Page \thepage \hspace{1pt} of \pageref{LastPage}}
\fancyfoot[R]{Page \thepage \hspace{1pt}  of \pageref{LastPage}}
% \fancyfoot[C]{
%     \footnotesize
%     \textcolor{customcolor}{\itshape\yourweb}\quad
%     \textcolor{customcolor}{\youremail}}
\fancyfoot[C]{%
  \footnotesize
  \emailsymbol\youremail \quad $\bullet$ \quad
 \itshape\homepagesymbol\yourweb
}
    
\renewcommand{\headrulewidth}{0pt}
\renewcommand{\footrulewidth}{0pt}

% First page carries the contact block in the title, so suppress the
% email/website footer there (keep only the page number).
\fancypagestyle{firstpage}{%
  \fancyhf{}%
  \fancyfoot[R]{Page \thepage \hspace{1pt}  of \pageref{LastPage}}%
  \renewcommand{\headrulewidth}{0pt}%
  \renewcommand{\footrulewidth}{0pt}%
}

\usepackage{lastpage}

%\lhead{This is my name}
%\rhead{this is page \thepage}
%\renewcommand{\headrulewidth}{0.4pt}
%\renewcommand{\footrulewidth}{0.4pt}

%debug page outer frames
%\usepackage{showframe}

%-----------------------
%  commands for circles 
%----------------------------------------------------------------------------------------
%	BEGIN DOCUMENT
%----------------------------------------------------------------------------------------
%\pagenumbering{arabic}

% ---------------------------------------------------------------------------
% Per-application tailoring. Wrapper files (e.g., Kline_CV_washu.tex) define
% these BEFORE %-----------------------------------------------------------------------------------------------------------------------------------------------%
%	The MIT License (MIT)
%
%	Copyright (c) 2021 Jitin Nair
%
%	Permission is hereby granted, free of charge, to any person obtaining a copy
%	of this software and associated documentation files (the "Software"), to deal
%	in the Software without restriction, including without limitation the rights
%	to use, copy, modify, merge, publish, distribute, sublicense, and/or sell
%	copies of the Software, and to permit persons to whom the Software is
%	furnished to do so, subject to the following conditions:
%	
%	THE SOFTWARE IS PROVIDED "AS IS", WITHOUT WARRANTY OF ANY KIND, EXPRESS OR
%	IMPLIED, INCLUDING BUT NOT LIMITED TO THE WARRANTIES OF MERCHANTABILITY,
%	FITNESS FOR A PARTICULAR PURPOSE AND NONINFRINGEMENT. IN NO EVENT SHALL THE
%	AUTHORS OR COPYRIGHT HOLDERS BE LIABLE FOR ANY CLAIM, DAMAGES OR OTHER
%	LIABILITY, WHETHER IN AN ACTION OF CONTRACT, TORT OR OTHERWISE, ARISING FROM,
%	OUT OF OR IN CONNECTION WITH THE SOFTWARE OR THE USE OR OTHER DEALINGS IN
%	THE SOFTWARE.
%	
%
%-----------------------------------------------------------------------------------------------------------------------------------------------%

%----------------------------------------------------------------------------------------
%	DOCUMENT DEFINITION
%----------------------------------------------------------------------------------------

% article class because we want to fully customize the page and not use a cv template
\documentclass[letterpaper,10pt]{article}

% Adding packages for skills
% -------------------------
\usepackage{geometry}
\geometry{margin=0.75in}

\usepackage{array}
\usepackage[svgnames]{xcolor}
\usepackage{tikz}
\usetikzlibrary{math}
\usepackage{progressbar}
%\usepackage[maxbibnames=99]{biblatex}
\usepackage{url}
%\usepackage[backend=biber,style=numeric,sorting=ydnt,defernumbers=true]{biblatex}
% Bibliography is typeset manually in the Publications / Talks sections below
% (rendered from publications.md); biblatex/biber is intentionally not used.
% \usepackage[style=numeric, sorting=none, backend=biber, maxbibnames=99]{biblatex}
% \addbibresource{citations.bib}


%----------------------------------------------------------------------------------------
%	FONT
%----------------------------------------------------------------------------------------

% Same font stack as the research and teaching statements:
% lmodern first gives T1 a scalable fallback so nothing lands on a bitmap
% Computer Modern face; 'default' is what actually makes Lato the document
% font (without it the package loads but is never used).
\usepackage{lmodern}
\usepackage[T1]{fontenc}
\usepackage[scaled=0.98,default]{lato}


% % fontspec allows you to use TTF/OTF fonts directly
% \usepackage{fontspec}
% \defaultfontfeatures{Ligatures=TeX}

% % modified for ShareLaTeX use
% \setmainfont[
% SmallCapsFont = Fontin-SmallCaps.otf,
% BoldFont = Fontin-Bold.otf,
% ItalicFont = Fontin-Italic.otf
% ]
% {Fontin.otf}

%----------------------------------------------------------------------------------------
%	PACKAGES
%----------------------------------------------------------------------------------------
\usepackage{url}
\usepackage{parskip}
\setlength{\parskip}{3.5pt plus 1pt}   % tighter inter-paragraph spacing than parskip default

%other packages for formatting
\RequirePackage{color}
\RequirePackage{graphicx}
\usepackage[usenames,dvipsnames]{xcolor}
%tabularx environment
%for lists within experience section
\usepackage{enumitem}

% centered version of 'X' col. type
\newcolumntype{C}{>{\centering\arraybackslash}X} 

%to prevent spillover of tabular into next pages
\usepackage{supertabular}
\usepackage{tabularx}
\usepackage{needspace}   % reserve vertical space so a block isn't split across pages
\newlength{\fullcollw}
\setlength{\fullcollw}{0.47\textwidth}
\newlength{\awardcolw}   % description column of the page-breakable awards table

%custom \section
\usepackage{titlesec}				
\usepackage{multicol}
\usepackage{multirow}

%CV Sections inspired by: 
%http://stefano.italians.nl/archives/26
\titleformat{\section}{\normalsize\scshape\raggedright}{}{0em}{}[\titlerule]
%\titleformat{\section}{\textbf{}\normalsize\scshape\raggedright}{}{0em}{}[\titlerule]
\titlespacing{\section}{0pt}{7pt}{5pt}

%for publications

%Setup hyperref package, and colours for links
\definecolor{customcolor}{HTML}{6f784d}
\colorlet{color1}{customcolor}
\usepackage[unicode, draft=false]{hyperref}
\definecolor{linkcolour}{rgb}{0,0.2,0.6}
\definecolor{darknavy}{HTML}{1F3864}   % dark navy for hyperlinks (prints legibly)
\hypersetup{colorlinks,breaklinks,urlcolor=darknavy,linkcolor=darknavy,citecolor=darknavy}
\usepackage{xurl}   % allow long DOIs/URLs to wrap across lines
\urlstyle{same}     % use the body font for wrapped identifiers, not typewriter
\hypersetup{pdftitle={Jenna M. Kline --- Curriculum Vitae}, pdfauthor={Jenna M. Kline}}

%--- Publication-list formatting (manual, rendered from publications.md) ---
\providecommand{\emailsymbol}{\faEnvelope~}
\providecommand{\homepagesymbol}{\faGlobe~}
% Compact small-caps subsection headers for publication categories
\titleformat{\subsection}{\normalsize\bfseries\raggedright}{}{0em}{}
\titlespacing{\subsection}{0pt}{5pt}{2pt}
% Hanging-indent list for references. Numbering restarts in every category and
% carries a category prefix (F/J/C/P/T/A) so cross-references stay unambiguous.
\newlist{pubs}{enumerate}{1}
\setlist[pubs]{leftmargin=2.5em, labelindent=0pt, labelwidth=2.1em, labelsep=0.4em, align=left, label=\arabic*., ref=\arabic*, itemsep=0.45em, topsep=0.12em, parsep=0pt}
% Year sub-heading within a category
\newcommand{\pubyear}[1]{\smallskip\noindent\textbf{#1}\par\nobreak\vspace{0.15em}}
% Clickable identifiers
\newcommand{\doilink}[1]{\href{https://doi.org/#1}{doi:\nolinkurl{#1}}}
% Artifact host link (GitHub, Hugging Face, ...). \textnormal keeps it upright
% when it sits inside the italic venue field of \pvenue.
\newcommand{\phost}[2]{\textnormal{\href{#1}{#2}}}

% --- Three-line publication entry: title / authors / venue + year -----------
% Title carries the emphasis; venue and year sit on their own line so a reader
% can scan where and when the work appeared without parsing the author list.
\newcommand{\ptitle}[1]{\textbf{#1}\par\nobreak\vspace{0.4pt}}
\newcommand{\pauth}[1]{#1\par\nobreak\vspace{0.4pt}}
% #1 = venue, #2 = year, #3 = identifiers, awards, and status notes (may be empty)
\newcommand{\pvenue}[3]{\textit{#1}, \textbf{#2}. #3\par}
% Artifact detail line: everything but the title and the download count runs on
% as one paragraph. #1 = kind, hosts, and role (upright); #2 = contents and the
% publications it supports (italic, like the venue field of \pvenue).
\newcommand{\pdetail}[2]{#1 \textit{#2}\par}
% Status flag for work that is not yet published
\newcommand{\status}[1]{\textbf{[#1]}}
% Category sub-heading inside a publication list block
\newcommand{\talkgroup}[1]{\smallskip\noindent\textit{\textbf{#1}}\par\nobreak\vspace{0.2em}}
% Impact metric on its own line, below the venue line of an artifact entry
\newcommand{\pmetric}[1]{\textbf{#1}\par}
% Qualifier trailing a subsection heading on the same line. \normalfont undoes
% the heading's bold; \footnotesize\itshape keeps it subordinate to the title.
\newcommand{\subnote}[1]{{\normalfont\footnotesize\itshape #1}}
% Field-experience entry: location (left) + dates (right) on line 1, description on line 2
% #1 = location, #2 = dates, #3 = description
\newcommand{\fieldentry}[3]{%
  \noindent\begin{tabularx}{\linewidth}{@{}X r@{}}\textbf{#1} & #2\end{tabularx}\par\nobreak\vspace{1pt}%
  {\noindent #3\par}\vspace{2pt}}

%for social icons
\usepackage{fontawesome5}
\usepackage{fontawesome5}
\usepackage{graphicx}  

% for bib
%\usepackage{cite}



% %for header
% \usepackage{fancyhdr}
% \pagestyle{fancy}
% \fancyhf{}
% \fancyfoot[R]{Page \thepage \hspace{1pt} of 4}
% \fancyfoot[L]{Jenna M. Kline}
% \renewcommand{\headrulewidth}{0pt}
% \renewcommand{\footrulewidth}{0pt}

%----------------------------------------------------------------------------------------
%	HEADER AND FOOTER DEFINITIONS
%----------------------------------------------------------------------------------------
\newcommand{\yourname}{Jenna M. Kline}
% \newcommand{\youremail}{jennamkline@gmail.com}
% \newcommand{\yourweb}{jennamk14.github.io}

\newcommand{\youremail}{\href{mailto:jennamk9@mit.edu}{jennamk9@mit.edu}}
\newcommand{\yourweb}{\href{https://jennamk14.github.io/}{jennamk14.github.io}}

%for header
\usepackage{fancyhdr}
\pagestyle{fancy}
\fancyhf{}
% \fancyhead[L]{\yourname\quad}
% \fancyhead[R]{Curriculum Vitae}
%\fancyfoot[R]{Page \thepage \hspace{1pt} of \pageref{LastPage}}
\fancyfoot[R]{Page \thepage \hspace{1pt}  of \pageref{LastPage}}
% \fancyfoot[C]{
%     \footnotesize
%     \textcolor{customcolor}{\itshape\yourweb}\quad
%     \textcolor{customcolor}{\youremail}}
\fancyfoot[C]{%
  \footnotesize
  \emailsymbol\youremail \quad $\bullet$ \quad
 \itshape\homepagesymbol\yourweb
}
    
\renewcommand{\headrulewidth}{0pt}
\renewcommand{\footrulewidth}{0pt}

% First page carries the contact block in the title, so suppress the
% email/website footer there (keep only the page number).
\fancypagestyle{firstpage}{%
  \fancyhf{}%
  \fancyfoot[R]{Page \thepage \hspace{1pt}  of \pageref{LastPage}}%
  \renewcommand{\headrulewidth}{0pt}%
  \renewcommand{\footrulewidth}{0pt}%
}

\usepackage{lastpage}

%\lhead{This is my name}
%\rhead{this is page \thepage}
%\renewcommand{\headrulewidth}{0.4pt}
%\renewcommand{\footrulewidth}{0.4pt}

%debug page outer frames
%\usepackage{showframe}

%-----------------------
%  commands for circles 
%----------------------------------------------------------------------------------------
%	BEGIN DOCUMENT
%----------------------------------------------------------------------------------------
%\pagenumbering{arabic}

% ---------------------------------------------------------------------------
% Per-application tailoring. Wrapper files (e.g., Kline_CV_washu.tex) define
% these BEFORE %-----------------------------------------------------------------------------------------------------------------------------------------------%
%	The MIT License (MIT)
%
%	Copyright (c) 2021 Jitin Nair
%
%	Permission is hereby granted, free of charge, to any person obtaining a copy
%	of this software and associated documentation files (the "Software"), to deal
%	in the Software without restriction, including without limitation the rights
%	to use, copy, modify, merge, publish, distribute, sublicense, and/or sell
%	copies of the Software, and to permit persons to whom the Software is
%	furnished to do so, subject to the following conditions:
%	
%	THE SOFTWARE IS PROVIDED "AS IS", WITHOUT WARRANTY OF ANY KIND, EXPRESS OR
%	IMPLIED, INCLUDING BUT NOT LIMITED TO THE WARRANTIES OF MERCHANTABILITY,
%	FITNESS FOR A PARTICULAR PURPOSE AND NONINFRINGEMENT. IN NO EVENT SHALL THE
%	AUTHORS OR COPYRIGHT HOLDERS BE LIABLE FOR ANY CLAIM, DAMAGES OR OTHER
%	LIABILITY, WHETHER IN AN ACTION OF CONTRACT, TORT OR OTHERWISE, ARISING FROM,
%	OUT OF OR IN CONNECTION WITH THE SOFTWARE OR THE USE OR OTHER DEALINGS IN
%	THE SOFTWARE.
%	
%
%-----------------------------------------------------------------------------------------------------------------------------------------------%

%----------------------------------------------------------------------------------------
%	DOCUMENT DEFINITION
%----------------------------------------------------------------------------------------

% article class because we want to fully customize the page and not use a cv template
\documentclass[letterpaper,10pt]{article}

% Adding packages for skills
% -------------------------
\usepackage{geometry}
\geometry{margin=0.75in}

\usepackage{array}
\usepackage[svgnames]{xcolor}
\usepackage{tikz}
\usetikzlibrary{math}
\usepackage{progressbar}
%\usepackage[maxbibnames=99]{biblatex}
\usepackage{url}
%\usepackage[backend=biber,style=numeric,sorting=ydnt,defernumbers=true]{biblatex}
% Bibliography is typeset manually in the Publications / Talks sections below
% (rendered from publications.md); biblatex/biber is intentionally not used.
% \usepackage[style=numeric, sorting=none, backend=biber, maxbibnames=99]{biblatex}
% \addbibresource{citations.bib}


%----------------------------------------------------------------------------------------
%	FONT
%----------------------------------------------------------------------------------------

% Same font stack as the research and teaching statements:
% lmodern first gives T1 a scalable fallback so nothing lands on a bitmap
% Computer Modern face; 'default' is what actually makes Lato the document
% font (without it the package loads but is never used).
\usepackage{lmodern}
\usepackage[T1]{fontenc}
\usepackage[scaled=0.98,default]{lato}


% % fontspec allows you to use TTF/OTF fonts directly
% \usepackage{fontspec}
% \defaultfontfeatures{Ligatures=TeX}

% % modified for ShareLaTeX use
% \setmainfont[
% SmallCapsFont = Fontin-SmallCaps.otf,
% BoldFont = Fontin-Bold.otf,
% ItalicFont = Fontin-Italic.otf
% ]
% {Fontin.otf}

%----------------------------------------------------------------------------------------
%	PACKAGES
%----------------------------------------------------------------------------------------
\usepackage{url}
\usepackage{parskip}
\setlength{\parskip}{3.5pt plus 1pt}   % tighter inter-paragraph spacing than parskip default

%other packages for formatting
\RequirePackage{color}
\RequirePackage{graphicx}
\usepackage[usenames,dvipsnames]{xcolor}
%tabularx environment
%for lists within experience section
\usepackage{enumitem}

% centered version of 'X' col. type
\newcolumntype{C}{>{\centering\arraybackslash}X} 

%to prevent spillover of tabular into next pages
\usepackage{supertabular}
\usepackage{tabularx}
\usepackage{needspace}   % reserve vertical space so a block isn't split across pages
\newlength{\fullcollw}
\setlength{\fullcollw}{0.47\textwidth}
\newlength{\awardcolw}   % description column of the page-breakable awards table

%custom \section
\usepackage{titlesec}				
\usepackage{multicol}
\usepackage{multirow}

%CV Sections inspired by: 
%http://stefano.italians.nl/archives/26
\titleformat{\section}{\normalsize\scshape\raggedright}{}{0em}{}[\titlerule]
%\titleformat{\section}{\textbf{}\normalsize\scshape\raggedright}{}{0em}{}[\titlerule]
\titlespacing{\section}{0pt}{7pt}{5pt}

%for publications

%Setup hyperref package, and colours for links
\definecolor{customcolor}{HTML}{6f784d}
\colorlet{color1}{customcolor}
\usepackage[unicode, draft=false]{hyperref}
\definecolor{linkcolour}{rgb}{0,0.2,0.6}
\definecolor{darknavy}{HTML}{1F3864}   % dark navy for hyperlinks (prints legibly)
\hypersetup{colorlinks,breaklinks,urlcolor=darknavy,linkcolor=darknavy,citecolor=darknavy}
\usepackage{xurl}   % allow long DOIs/URLs to wrap across lines
\urlstyle{same}     % use the body font for wrapped identifiers, not typewriter
\hypersetup{pdftitle={Jenna M. Kline --- Curriculum Vitae}, pdfauthor={Jenna M. Kline}}

%--- Publication-list formatting (manual, rendered from publications.md) ---
\providecommand{\emailsymbol}{\faEnvelope~}
\providecommand{\homepagesymbol}{\faGlobe~}
% Compact small-caps subsection headers for publication categories
\titleformat{\subsection}{\normalsize\bfseries\raggedright}{}{0em}{}
\titlespacing{\subsection}{0pt}{5pt}{2pt}
% Hanging-indent list for references. Numbering restarts in every category and
% carries a category prefix (F/J/C/P/T/A) so cross-references stay unambiguous.
\newlist{pubs}{enumerate}{1}
\setlist[pubs]{leftmargin=2.5em, labelindent=0pt, labelwidth=2.1em, labelsep=0.4em, align=left, label=\arabic*., ref=\arabic*, itemsep=0.45em, topsep=0.12em, parsep=0pt}
% Year sub-heading within a category
\newcommand{\pubyear}[1]{\smallskip\noindent\textbf{#1}\par\nobreak\vspace{0.15em}}
% Clickable identifiers
\newcommand{\doilink}[1]{\href{https://doi.org/#1}{doi:\nolinkurl{#1}}}
% Artifact host link (GitHub, Hugging Face, ...). \textnormal keeps it upright
% when it sits inside the italic venue field of \pvenue.
\newcommand{\phost}[2]{\textnormal{\href{#1}{#2}}}

% --- Three-line publication entry: title / authors / venue + year -----------
% Title carries the emphasis; venue and year sit on their own line so a reader
% can scan where and when the work appeared without parsing the author list.
\newcommand{\ptitle}[1]{\textbf{#1}\par\nobreak\vspace{0.4pt}}
\newcommand{\pauth}[1]{#1\par\nobreak\vspace{0.4pt}}
% #1 = venue, #2 = year, #3 = identifiers, awards, and status notes (may be empty)
\newcommand{\pvenue}[3]{\textit{#1}, \textbf{#2}. #3\par}
% Artifact detail line: everything but the title and the download count runs on
% as one paragraph. #1 = kind, hosts, and role (upright); #2 = contents and the
% publications it supports (italic, like the venue field of \pvenue).
\newcommand{\pdetail}[2]{#1 \textit{#2}\par}
% Status flag for work that is not yet published
\newcommand{\status}[1]{\textbf{[#1]}}
% Category sub-heading inside a publication list block
\newcommand{\talkgroup}[1]{\smallskip\noindent\textit{\textbf{#1}}\par\nobreak\vspace{0.2em}}
% Impact metric on its own line, below the venue line of an artifact entry
\newcommand{\pmetric}[1]{\textbf{#1}\par}
% Qualifier trailing a subsection heading on the same line. \normalfont undoes
% the heading's bold; \footnotesize\itshape keeps it subordinate to the title.
\newcommand{\subnote}[1]{{\normalfont\footnotesize\itshape #1}}
% Field-experience entry: location (left) + dates (right) on line 1, description on line 2
% #1 = location, #2 = dates, #3 = description
\newcommand{\fieldentry}[3]{%
  \noindent\begin{tabularx}{\linewidth}{@{}X r@{}}\textbf{#1} & #2\end{tabularx}\par\nobreak\vspace{1pt}%
  {\noindent #3\par}\vspace{2pt}}

%for social icons
\usepackage{fontawesome5}
\usepackage{fontawesome5}
\usepackage{graphicx}  

% for bib
%\usepackage{cite}



% %for header
% \usepackage{fancyhdr}
% \pagestyle{fancy}
% \fancyhf{}
% \fancyfoot[R]{Page \thepage \hspace{1pt} of 4}
% \fancyfoot[L]{Jenna M. Kline}
% \renewcommand{\headrulewidth}{0pt}
% \renewcommand{\footrulewidth}{0pt}

%----------------------------------------------------------------------------------------
%	HEADER AND FOOTER DEFINITIONS
%----------------------------------------------------------------------------------------
\newcommand{\yourname}{Jenna M. Kline}
% \newcommand{\youremail}{jennamkline@gmail.com}
% \newcommand{\yourweb}{jennamk14.github.io}

\newcommand{\youremail}{\href{mailto:jennamk9@mit.edu}{jennamk9@mit.edu}}
\newcommand{\yourweb}{\href{https://jennamk14.github.io/}{jennamk14.github.io}}

%for header
\usepackage{fancyhdr}
\pagestyle{fancy}
\fancyhf{}
% \fancyhead[L]{\yourname\quad}
% \fancyhead[R]{Curriculum Vitae}
%\fancyfoot[R]{Page \thepage \hspace{1pt} of \pageref{LastPage}}
\fancyfoot[R]{Page \thepage \hspace{1pt}  of \pageref{LastPage}}
% \fancyfoot[C]{
%     \footnotesize
%     \textcolor{customcolor}{\itshape\yourweb}\quad
%     \textcolor{customcolor}{\youremail}}
\fancyfoot[C]{%
  \footnotesize
  \emailsymbol\youremail \quad $\bullet$ \quad
 \itshape\homepagesymbol\yourweb
}
    
\renewcommand{\headrulewidth}{0pt}
\renewcommand{\footrulewidth}{0pt}

% First page carries the contact block in the title, so suppress the
% email/website footer there (keep only the page number).
\fancypagestyle{firstpage}{%
  \fancyhf{}%
  \fancyfoot[R]{Page \thepage \hspace{1pt}  of \pageref{LastPage}}%
  \renewcommand{\headrulewidth}{0pt}%
  \renewcommand{\footrulewidth}{0pt}%
}

\usepackage{lastpage}

%\lhead{This is my name}
%\rhead{this is page \thepage}
%\renewcommand{\headrulewidth}{0.4pt}
%\renewcommand{\footrulewidth}{0.4pt}

%debug page outer frames
%\usepackage{showframe}

%-----------------------
%  commands for circles 
%----------------------------------------------------------------------------------------
%	BEGIN DOCUMENT
%----------------------------------------------------------------------------------------
%\pagenumbering{arabic}

% ---------------------------------------------------------------------------
% Per-application tailoring. Wrapper files (e.g., Kline_CV_washu.tex) define
% these BEFORE \input{Kline_CV}; otherwise the defaults below are used.
%   \cvfocus: r1 (research first) or teaching (mentoring & teaching first)
% ---------------------------------------------------------------------------
\providecommand*{\cvfocus}{r1}
% \cvroot: path from the compiling directory to this CV folder (empty when
% compiling here). Private per-application wrappers set it, e.g. ../../../../CV/
\providecommand{\cvroot}{}
% A local ./sections/<name>.tex (next to a wrapper) overrides the shared section.
\newcommand{\cvsection}[1]{\IfFileExists{./sections/#1.tex}{\input{./sections/#1}}{\input{\cvroot sections/#1}}}
\def\cvteachingfocus{teaching}
\providecommand{\cvsummarytitle}{Cyber-Physical Systems for Adaptive Ecological Sensing at the Far Edge}
\providecommand{\cvsummary}{I build \textbf{cyber-physical systems} for autonomous environmental sensing in dynamic, resource-constrained field settings. My research advances \textbf{distributed intelligence}, \textbf{edge AI}, and \textbf{adaptive autonomy} to coordinate heterogeneous sensors and robots, enabling \textbf{cross-modal systems} to adapt what, where, and when they observe. I deploy these systems across \textbf{ecology}, \textbf{biodiversity monitoring}, and \textbf{coastal and climate resilience}, developing the foundations for persistent, adaptive observation of complex environmental systems.}

\begin{document}
\thispagestyle{firstpage}% page 1 shows only the page number in the footer

% non-numbered pages
%\pagestyle{empty}

%----------------------------------------------------------------------------------------
%	TITLE
%----------------------------------------------------------------------------------------

% \begin{tabularx}{\linewidth}{@{}X X@{} }
% \huge{Your Name}\vspace{2pt} & \hfill \emoji{incoming-envelope} email@email.com \\
% \raisebox{-0.05\height}\faGithub\ username \ | \
% \raisebox{-0.00\height}\faLinkedin\ username \ | \ \raisebox{-0.05\height}\faGlobe \ mysite.com  & \hfill \emoji{calling} number
% \end{tabularx}

\begin{tabularx}{\linewidth}{@{}C@{}}
\Large\textbf{{Jenna M. Kline}} \\[7.5pt]
\href{https://jennamk14.github.io/}{\raisebox{-0.05\height}\faGithub\ Github} \ $|$ \ 
%
\href{https://www.linkedin.com/in/jenna-kline-00937ab2/}{\raisebox{-0.05\height}\faLinkedin\ LinkedIn} \ $|$
%
\ 
%
\href{https://scholar.google.com/citations?user=bhNDuigAAAAJ}{\raisebox{-0.05\height}\faGraduationCap \ Google Scholar} \ $|$ \
%
\href{mailto:jennamk9@mit.edu}{\raisebox{-0.05\height}\faEnvelope \ jennamk9@mit.edu} \
%
%\href{mailto:jennamkline@gmail.com}{\raisebox{-0.05\height}\faEnvelope \ jennamkline@gmail.com} \
%
%
%\href{https://resume.io/r/m5UqgujdL}{\raisebox{-0.05\height}\faFile \ resume.io}\ 
%\href{tel:+00000000000}{\raisebox{-0.05\height}\faMobile \ +00.00.000.000} \\

\end{tabularx}

%----------------------------------------------------------------------------------------
%	EDUCATION
%----------------------------------------------------------------------------------------
\vspace{10pt}   % shrink whitespace between contact info and Education

% ---------------------------------------------------------------------------
% Section order (CV feedback, Sep 2026): tailored research summary first, then
% academics (education, appointments), then research for R1 applications or
% mentoring & teaching for undergraduate-focused applications.
% ---------------------------------------------------------------------------
\cvsection{summary}
\cvsection{education}
\cvsection{experience}
\ifx\cvfocus\cvteachingfocus
  \cvsection{mentoring}
  \cvsection{service}
  % featured publications appear inside research outputs in this order
  \cvsection{awards}
  \cvsection{field}
  \cvsection{skills}
\else
  \cvsection{featured}
  \textit{Remaining research outputs (journals, conferences, presentations, artifacts) included in the research outputs section (p.~\pageref{sec:research-outputs}).}
  \cvsection{awards}
  \cvsection{field}
  \cvsection{skills}
  \cvsection{mentoring}
  \cvsection{service}
\fi
\cvsection{outputs}
% \cvsection{references}

\end{document}


% ---------------------------------------------------------------------------
% ESA CV REVIEW FEEDBACK --- status
% ---------------------------------------------------------------------------
% 1. DONE - References section added at end (Tanya, Chris, Daniel Rubenstein,
%    Tilo Burghardt). TODO: fill in email addresses (marked in red).
% 2. DONE - Skills section added as scannable keyword lists (field research,
%    drones/sensing, edge AI, data infrastructure, HPC, programming). Placed
%    after Featured Publications per the section order in item 6.
% 3. DONE - Numbering restarts in each category with a category prefix
%    (F/J/C/P/A) so cross-references stay unambiguous. Entries reformatted to
%    title / authors / venue+year, with the venue italicised and the year bold.
%    Preprints redistributed: FAIR2 Drones + cross-modal blueprint -> Featured;
%    "Where to Trust", Pollock et al., Meier et al., Aguzzi et al. -> Journals;
%    WildBox -> Conference; kabr-tools -> Artifacts. All marked [Under review].
%    The two ACSOS 2026 papers sit under Conference, not Featured.
%    DOIs, page numbers, and arXiv links are omitted throughout.
% 4. DONE - All-time Hugging Face download totals added to the three dataset
%    collections (KABR 164K, MMLA 204K, SmartWilds 13K, as of September 2026).
%    These need refreshing before each submission.
% 5. DONE - Talks listed as a single reverse-chronological list (the
%    group-by-community subheadings were removed).
% 6. DONE - Section order: education, research interests, academic appointments,
%    featured publications, skills, grants and awards, field experience,
%    mentoring and teaching, service, research outputs (journals, conferences,
%    presentations, artifacts), references.
% 7. DONE - "Student Supervision and Teaching Experience" renamed to
%    "Mentoring and Teaching", with mentoring listed first.
%
% 8. DONE - WildDroneEU entry expanded (ESA feedback): now states that the work
%    was independent dissertation research within the consortium, and itemizes
%    the two research visits, the summer school, and the two Kenya field
%    campaigns (own systems tested + support for collaborators).

; otherwise the defaults below are used.
%   \cvfocus: r1 (research first) or teaching (mentoring & teaching first)
% ---------------------------------------------------------------------------
\providecommand*{\cvfocus}{r1}
% \cvroot: path from the compiling directory to this CV folder (empty when
% compiling here). Private per-application wrappers set it, e.g. ../../../../CV/
\providecommand{\cvroot}{}
% A local ./sections/<name>.tex (next to a wrapper) overrides the shared section.
\newcommand{\cvsection}[1]{\IfFileExists{./sections/#1.tex}{\input{./sections/#1}}{\input{\cvroot sections/#1}}}
\def\cvteachingfocus{teaching}
\providecommand{\cvsummarytitle}{Cyber-Physical Systems for Adaptive Ecological Sensing at the Far Edge}
\providecommand{\cvsummary}{I build \textbf{cyber-physical systems} for autonomous environmental sensing in dynamic, resource-constrained field settings. My research advances \textbf{distributed intelligence}, \textbf{edge AI}, and \textbf{adaptive autonomy} to coordinate heterogeneous sensors and robots, enabling \textbf{cross-modal systems} to adapt what, where, and when they observe. I deploy these systems across \textbf{ecology}, \textbf{biodiversity monitoring}, and \textbf{coastal and climate resilience}, developing the foundations for persistent, adaptive observation of complex environmental systems.}

\begin{document}
\thispagestyle{firstpage}% page 1 shows only the page number in the footer

% non-numbered pages
%\pagestyle{empty}

%----------------------------------------------------------------------------------------
%	TITLE
%----------------------------------------------------------------------------------------

% \begin{tabularx}{\linewidth}{@{}X X@{} }
% \huge{Your Name}\vspace{2pt} & \hfill \emoji{incoming-envelope} email@email.com \\
% \raisebox{-0.05\height}\faGithub\ username \ | \
% \raisebox{-0.00\height}\faLinkedin\ username \ | \ \raisebox{-0.05\height}\faGlobe \ mysite.com  & \hfill \emoji{calling} number
% \end{tabularx}

\begin{tabularx}{\linewidth}{@{}C@{}}
\Large\textbf{{Jenna M. Kline}} \\[7.5pt]
\href{https://jennamk14.github.io/}{\raisebox{-0.05\height}\faGithub\ Github} \ $|$ \ 
%
\href{https://www.linkedin.com/in/jenna-kline-00937ab2/}{\raisebox{-0.05\height}\faLinkedin\ LinkedIn} \ $|$
%
\ 
%
\href{https://scholar.google.com/citations?user=bhNDuigAAAAJ}{\raisebox{-0.05\height}\faGraduationCap \ Google Scholar} \ $|$ \
%
\href{mailto:jennamk9@mit.edu}{\raisebox{-0.05\height}\faEnvelope \ jennamk9@mit.edu} \
%
%\href{mailto:jennamkline@gmail.com}{\raisebox{-0.05\height}\faEnvelope \ jennamkline@gmail.com} \
%
%
%\href{https://resume.io/r/m5UqgujdL}{\raisebox{-0.05\height}\faFile \ resume.io}\ 
%\href{tel:+00000000000}{\raisebox{-0.05\height}\faMobile \ +00.00.000.000} \\

\end{tabularx}

%----------------------------------------------------------------------------------------
%	EDUCATION
%----------------------------------------------------------------------------------------
\vspace{10pt}   % shrink whitespace between contact info and Education

% ---------------------------------------------------------------------------
% Section order (CV feedback, Sep 2026): tailored research summary first, then
% academics (education, appointments), then research for R1 applications or
% mentoring & teaching for undergraduate-focused applications.
% ---------------------------------------------------------------------------
\cvsection{summary}
\cvsection{education}
\cvsection{experience}
\ifx\cvfocus\cvteachingfocus
  \cvsection{mentoring}
  \cvsection{service}
  % featured publications appear inside research outputs in this order
  \cvsection{awards}
  \cvsection{field}
  \cvsection{skills}
\else
  \cvsection{featured}
  \textit{Remaining research outputs (journals, conferences, presentations, artifacts) included in the research outputs section (p.~\pageref{sec:research-outputs}).}
  \cvsection{awards}
  \cvsection{field}
  \cvsection{skills}
  \cvsection{mentoring}
  \cvsection{service}
\fi
\cvsection{outputs}
% \cvsection{references}

\end{document}


% ---------------------------------------------------------------------------
% ESA CV REVIEW FEEDBACK --- status
% ---------------------------------------------------------------------------
% 1. DONE - References section added at end (Tanya, Chris, Daniel Rubenstein,
%    Tilo Burghardt). TODO: fill in email addresses (marked in red).
% 2. DONE - Skills section added as scannable keyword lists (field research,
%    drones/sensing, edge AI, data infrastructure, HPC, programming). Placed
%    after Featured Publications per the section order in item 6.
% 3. DONE - Numbering restarts in each category with a category prefix
%    (F/J/C/P/A) so cross-references stay unambiguous. Entries reformatted to
%    title / authors / venue+year, with the venue italicised and the year bold.
%    Preprints redistributed: FAIR2 Drones + cross-modal blueprint -> Featured;
%    "Where to Trust", Pollock et al., Meier et al., Aguzzi et al. -> Journals;
%    WildBox -> Conference; kabr-tools -> Artifacts. All marked [Under review].
%    The two ACSOS 2026 papers sit under Conference, not Featured.
%    DOIs, page numbers, and arXiv links are omitted throughout.
% 4. DONE - All-time Hugging Face download totals added to the three dataset
%    collections (KABR 164K, MMLA 204K, SmartWilds 13K, as of September 2026).
%    These need refreshing before each submission.
% 5. DONE - Talks listed as a single reverse-chronological list (the
%    group-by-community subheadings were removed).
% 6. DONE - Section order: education, research interests, academic appointments,
%    featured publications, skills, grants and awards, field experience,
%    mentoring and teaching, service, research outputs (journals, conferences,
%    presentations, artifacts), references.
% 7. DONE - "Student Supervision and Teaching Experience" renamed to
%    "Mentoring and Teaching", with mentoring listed first.
%
% 8. DONE - WildDroneEU entry expanded (ESA feedback): now states that the work
%    was independent dissertation research within the consortium, and itemizes
%    the two research visits, the summer school, and the two Kenya field
%    campaigns (own systems tested + support for collaborators).

; otherwise the defaults below are used.
%   \cvfocus: r1 (research first) or teaching (mentoring & teaching first)
% ---------------------------------------------------------------------------
\providecommand*{\cvfocus}{r1}
% \cvroot: path from the compiling directory to this CV folder (empty when
% compiling here). Private per-application wrappers set it, e.g. ../../../../CV/
\providecommand{\cvroot}{}
% A local ./sections/<name>.tex (next to a wrapper) overrides the shared section.
\newcommand{\cvsection}[1]{\IfFileExists{./sections/#1.tex}{\input{./sections/#1}}{\input{\cvroot sections/#1}}}
\def\cvteachingfocus{teaching}
\providecommand{\cvsummarytitle}{Cyber-Physical Systems for Adaptive Ecological Sensing at the Far Edge}
\providecommand{\cvsummary}{I build \textbf{cyber-physical systems} for autonomous environmental sensing in dynamic, resource-constrained field settings. My research advances \textbf{distributed intelligence}, \textbf{edge AI}, and \textbf{adaptive autonomy} to coordinate heterogeneous sensors and robots, enabling \textbf{cross-modal systems} to adapt what, where, and when they observe. I deploy these systems across \textbf{ecology}, \textbf{biodiversity monitoring}, and \textbf{coastal and climate resilience}, developing the foundations for persistent, adaptive observation of complex environmental systems.}

\begin{document}
\thispagestyle{firstpage}% page 1 shows only the page number in the footer

% non-numbered pages
%\pagestyle{empty}

%----------------------------------------------------------------------------------------
%	TITLE
%----------------------------------------------------------------------------------------

% \begin{tabularx}{\linewidth}{@{}X X@{} }
% \huge{Your Name}\vspace{2pt} & \hfill \emoji{incoming-envelope} email@email.com \\
% \raisebox{-0.05\height}\faGithub\ username \ | \
% \raisebox{-0.00\height}\faLinkedin\ username \ | \ \raisebox{-0.05\height}\faGlobe \ mysite.com  & \hfill \emoji{calling} number
% \end{tabularx}

\begin{tabularx}{\linewidth}{@{}C@{}}
\Large\textbf{{Jenna M. Kline}} \\[7.5pt]
\href{https://jennamk14.github.io/}{\raisebox{-0.05\height}\faGithub\ Github} \ $|$ \ 
%
\href{https://www.linkedin.com/in/jenna-kline-00937ab2/}{\raisebox{-0.05\height}\faLinkedin\ LinkedIn} \ $|$
%
\ 
%
\href{https://scholar.google.com/citations?user=bhNDuigAAAAJ}{\raisebox{-0.05\height}\faGraduationCap \ Google Scholar} \ $|$ \
%
\href{mailto:jennamk9@mit.edu}{\raisebox{-0.05\height}\faEnvelope \ jennamk9@mit.edu} \
%
%\href{mailto:jennamkline@gmail.com}{\raisebox{-0.05\height}\faEnvelope \ jennamkline@gmail.com} \
%
%
%\href{https://resume.io/r/m5UqgujdL}{\raisebox{-0.05\height}\faFile \ resume.io}\ 
%\href{tel:+00000000000}{\raisebox{-0.05\height}\faMobile \ +00.00.000.000} \\

\end{tabularx}

%----------------------------------------------------------------------------------------
%	EDUCATION
%----------------------------------------------------------------------------------------
\vspace{10pt}   % shrink whitespace between contact info and Education

% ---------------------------------------------------------------------------
% Section order (CV feedback, Sep 2026): tailored research summary first, then
% academics (education, appointments), then research for R1 applications or
% mentoring & teaching for undergraduate-focused applications.
% ---------------------------------------------------------------------------
\cvsection{summary}
\cvsection{education}
\cvsection{experience}
\ifx\cvfocus\cvteachingfocus
  \cvsection{mentoring}
  \cvsection{service}
  % featured publications appear inside research outputs in this order
  \cvsection{awards}
  \cvsection{field}
  \cvsection{skills}
\else
  \cvsection{featured}
  \textit{Remaining research outputs (journals, conferences, presentations, artifacts) included in the research outputs section (p.~\pageref{sec:research-outputs}).}
  \cvsection{awards}
  \cvsection{field}
  \cvsection{skills}
  \cvsection{mentoring}
  \cvsection{service}
\fi
\cvsection{outputs}
% \cvsection{references}

\end{document}


% ---------------------------------------------------------------------------
% ESA CV REVIEW FEEDBACK --- status
% ---------------------------------------------------------------------------
% 1. DONE - References section added at end (Tanya, Chris, Daniel Rubenstein,
%    Tilo Burghardt). TODO: fill in email addresses (marked in red).
% 2. DONE - Skills section added as scannable keyword lists (field research,
%    drones/sensing, edge AI, data infrastructure, HPC, programming). Placed
%    after Featured Publications per the section order in item 6.
% 3. DONE - Numbering restarts in each category with a category prefix
%    (F/J/C/P/A) so cross-references stay unambiguous. Entries reformatted to
%    title / authors / venue+year, with the venue italicised and the year bold.
%    Preprints redistributed: FAIR2 Drones + cross-modal blueprint -> Featured;
%    "Where to Trust", Pollock et al., Meier et al., Aguzzi et al. -> Journals;
%    WildBox -> Conference; kabr-tools -> Artifacts. All marked [Under review].
%    The two ACSOS 2026 papers sit under Conference, not Featured.
%    DOIs, page numbers, and arXiv links are omitted throughout.
% 4. DONE - All-time Hugging Face download totals added to the three dataset
%    collections (KABR 164K, MMLA 204K, SmartWilds 13K, as of September 2026).
%    These need refreshing before each submission.
% 5. DONE - Talks listed as a single reverse-chronological list (the
%    group-by-community subheadings were removed).
% 6. DONE - Section order: education, research interests, academic appointments,
%    featured publications, skills, grants and awards, field experience,
%    mentoring and teaching, service, research outputs (journals, conferences,
%    presentations, artifacts), references.
% 7. DONE - "Student Supervision and Teaching Experience" renamed to
%    "Mentoring and Teaching", with mentoring listed first.
%
% 8. DONE - WildDroneEU entry expanded (ESA feedback): now states that the work
%    was independent dissertation research within the consortium, and itemizes
%    the two research visits, the summer school, and the two Kenya field
%    campaigns (own systems tested + support for collaborators).

; otherwise the defaults below are used.
%   \cvfocus: r1 (research first) or teaching (mentoring & teaching first)
% ---------------------------------------------------------------------------
\providecommand*{\cvfocus}{r1}
% \cvroot: path from the compiling directory to this CV folder (empty when
% compiling here). Private per-application wrappers set it, e.g. ../../../../CV/
\providecommand{\cvroot}{}
% A local ./sections/<name>.tex (next to a wrapper) overrides the shared section.
\newcommand{\cvsection}[1]{\IfFileExists{./sections/#1.tex}{\input{./sections/#1}}{\input{\cvroot sections/#1}}}
\def\cvteachingfocus{teaching}
\providecommand{\cvsummarytitle}{Cyber-Physical Systems for Adaptive Ecological Sensing at the Far Edge}
\providecommand{\cvsummary}{I build \textbf{cyber-physical systems} for autonomous environmental sensing in dynamic, resource-constrained field settings. My research advances \textbf{distributed intelligence}, \textbf{edge AI}, and \textbf{adaptive autonomy} to coordinate heterogeneous sensors and robots, enabling \textbf{cross-modal systems} to adapt what, where, and when they observe. I deploy these systems across \textbf{ecology}, \textbf{biodiversity monitoring}, and \textbf{coastal and climate resilience}, developing the foundations for persistent, adaptive observation of complex environmental systems.}

\begin{document}
\thispagestyle{firstpage}% page 1 shows only the page number in the footer

% non-numbered pages
%\pagestyle{empty}

%----------------------------------------------------------------------------------------
%	TITLE
%----------------------------------------------------------------------------------------

% \begin{tabularx}{\linewidth}{@{}X X@{} }
% \huge{Your Name}\vspace{2pt} & \hfill \emoji{incoming-envelope} email@email.com \\
% \raisebox{-0.05\height}\faGithub\ username \ | \
% \raisebox{-0.00\height}\faLinkedin\ username \ | \ \raisebox{-0.05\height}\faGlobe \ mysite.com  & \hfill \emoji{calling} number
% \end{tabularx}

\begin{tabularx}{\linewidth}{@{}C@{}}
\Large\textbf{{Jenna M. Kline}} \\[7.5pt]
\href{https://jennamk14.github.io/}{\raisebox{-0.05\height}\faGithub\ Github} \ $|$ \ 
%
\href{https://www.linkedin.com/in/jenna-kline-00937ab2/}{\raisebox{-0.05\height}\faLinkedin\ LinkedIn} \ $|$
%
\ 
%
\href{https://scholar.google.com/citations?user=bhNDuigAAAAJ}{\raisebox{-0.05\height}\faGraduationCap \ Google Scholar} \ $|$ \
%
\href{mailto:jennamk9@mit.edu}{\raisebox{-0.05\height}\faEnvelope \ jennamk9@mit.edu} \
%
%\href{mailto:jennamkline@gmail.com}{\raisebox{-0.05\height}\faEnvelope \ jennamkline@gmail.com} \
%
%
%\href{https://resume.io/r/m5UqgujdL}{\raisebox{-0.05\height}\faFile \ resume.io}\ 
%\href{tel:+00000000000}{\raisebox{-0.05\height}\faMobile \ +00.00.000.000} \\

\end{tabularx}

%----------------------------------------------------------------------------------------
%	EDUCATION
%----------------------------------------------------------------------------------------
\vspace{10pt}   % shrink whitespace between contact info and Education

% ---------------------------------------------------------------------------
% Section order (CV feedback, Sep 2026): tailored research summary first, then
% academics (education, appointments), then research for R1 applications or
% mentoring & teaching for undergraduate-focused applications.
% ---------------------------------------------------------------------------
\cvsection{summary}
\cvsection{education}
\cvsection{experience}
\ifx\cvfocus\cvteachingfocus
  \cvsection{mentoring}
  \cvsection{service}
  % featured publications appear inside research outputs in this order
  \cvsection{awards}
  \cvsection{field}
  \cvsection{skills}
\else
  \cvsection{featured}
  \textit{Remaining research outputs (journals, conferences, presentations, artifacts) included in the research outputs section (p.~\pageref{sec:research-outputs}).}
  \cvsection{awards}
  \cvsection{field}
  \cvsection{skills}
  \cvsection{mentoring}
  \cvsection{service}
\fi
\cvsection{outputs}
% \cvsection{references}

\end{document}


% ---------------------------------------------------------------------------
% ESA CV REVIEW FEEDBACK --- status
% ---------------------------------------------------------------------------
% 1. DONE - References section added at end (Tanya, Chris, Daniel Rubenstein,
%    Tilo Burghardt). TODO: fill in email addresses (marked in red).
% 2. DONE - Skills section added as scannable keyword lists (field research,
%    drones/sensing, edge AI, data infrastructure, HPC, programming). Placed
%    after Featured Publications per the section order in item 6.
% 3. DONE - Numbering restarts in each category with a category prefix
%    (F/J/C/P/A) so cross-references stay unambiguous. Entries reformatted to
%    title / authors / venue+year, with the venue italicised and the year bold.
%    Preprints redistributed: FAIR2 Drones + cross-modal blueprint -> Featured;
%    "Where to Trust", Pollock et al., Meier et al., Aguzzi et al. -> Journals;
%    WildBox -> Conference; kabr-tools -> Artifacts. All marked [Under review].
%    The two ACSOS 2026 papers sit under Conference, not Featured.
%    DOIs, page numbers, and arXiv links are omitted throughout.
% 4. DONE - All-time Hugging Face download totals added to the three dataset
%    collections (KABR 164K, MMLA 204K, SmartWilds 13K, as of September 2026).
%    These need refreshing before each submission.
% 5. DONE - Talks listed as a single reverse-chronological list (the
%    group-by-community subheadings were removed).
% 6. DONE - Section order: education, research interests, academic appointments,
%    featured publications, skills, grants and awards, field experience,
%    mentoring and teaching, service, research outputs (journals, conferences,
%    presentations, artifacts), references.
% 7. DONE - "Student Supervision and Teaching Experience" renamed to
%    "Mentoring and Teaching", with mentoring listed first.
%
% 8. DONE - WildDroneEU entry expanded (ESA feedback): now states that the work
%    was independent dissertation research within the consortium, and itemizes
%    the two research visits, the summer school, and the two Kenya field
%    campaigns (own systems tested + support for collaborators).

